\chapter{與苦主的衝突處理}

和人妻約會的時候，有一定可能性會進入與苦主的衝突之中。本章將詳細分析各種可能出現的衝突狀況和相應的應對方式。

\section{事前的衝突避免}

《黃帝內經》有言：「聖人不治已病治未病，不治已亂治未亂」。和治病救人一樣，最好的衝突解決方法，就是從一開始避免衝突。

避免衝突最重要的一點，就是避免去苦主家和人妻約會。有一次郭菊陽發信息給我，說她老公不在家，叫我去她自稱價值6000萬的波託菲諾純水岸的豪宅約會。我當時堅決拒絕了。只讓她來我在白石洲的出租屋約會。

無論做什麼事情，都是講究一個天時地利人和。在苦主家進行約會，一旦遇到突發情況，即使有「老公不在家」的天時，地利、人和都不在你這一邊。在酒店、自己的房子，你都有門禁可以隨時阻止外人進入。在苦主家你就沒有門禁了。設想一下，如果你和人妻在性愛的過程中，苦主進入了你們所在的臥室，觸發（物理）戰鬥。這個時候，你沒有任何護具、沒有任何盾牌、也不熟悉任何可以保護自己的地形。而苦主卻是在熟悉的戰場戰鬥。在戰鬥之中很容易處於下風。而如果你在美國等有城堡原則（Castle Doctrine）的地區，苦主有權利在自家不需要提前警告、撤退的時候合法使用武力。因此在這些地區更應該完全避開在苦主家約會。

所以最佳的方法、是盡量把人妻邀請到酒店、出租屋等場所進行約會。如果去酒店，記得一個人先開房，這樣避免兩個人同時出現在公安局的資料庫，這個資料庫很容易泄漏。有條件的話，也可以在工作日白天時候到辦公室約會。我之前工作過的亞馬遜的Twitch，還有我很多朋友所在的谷歌，辦公室都配備有共享午休室（nap room），還有免費的飲料、餐食。工作日的中午，一條微信把人妻約到公司，在午休室約會。完事之後，去公司餐廳拿一杯免費的熱拿鐵，再喫一碗東京粿\footnote{「東京粿」是潮汕話一直以來對Pho也就是東京（河內）的粿條的稱呼，潮汕以外一般翻譯爲「越南粉」。}（Pho）或者咖喱雞飯（Chicken Tikka Masala），又可以能量滿滿地訓練AI模型去了。這不比不中不洋的波託菲諾純水岸愜意一百倍？

室外的約會，也盡量避開苦主的活動範圍。比如郭菊陽的老公經常去深圳南山區的方大城，於是我就約她去福田區的Coco Park和平安金融中心商場。又比如春節的時候郭菊陽老公在汕頭潮南區，那麼我就約她去龍湖區喫莊氏祥記，去金平區喫達濠炒飯。

至於自己的住所，並不建議作爲約會地點。主要是避免人妻把自己當作女主人了，和避免分手後的爭吵問題。

如果你實在是沒有條件找到她家之外的約會場所（比如你不幸已婚、老婆是家庭主婦且工資卡上交），那麼也需要做好準備。就像1661年鄭成功攻佔臺灣的時候，要趁着漲潮進攻一樣。你要和海軍統領觀察潮汐那樣觀察苦主的生活節律，在苦主在家出現機率最低的那個時間與人妻約會。

具體來說，假設你和人妻約會的時間長度是$\delta$，然後在觀察建模之後，苦主在某個時刻$x$回家的機率爲$p_x$，那麼你需要找到一個時刻$t$，以最小化$$\int_{t}^{t+\delta}p_x \,dx$$

\section{在苦主家被發現的應對}

如果不幸在人妻家偶遇苦主，千萬不能慌張，需要沉着應對。

首先是要在遇到苦主之前，就熟讀本章，然後把可能出現的情況和應對措施先在腦海裏面過一遍。做好所有的準備。這樣子在真實面對苦主的時候就能胸有成竹，從容不迫。

第二步，就是要仔細閱讀前面與人妻交往的法理和道德分析。必須深刻理解以下三點：
\begin{enumerate}
\item 在新中國婦女解放之後，婦女有着100\%的性自主權，婚姻存在與否，都不能阻止女性和自己喜歡的男性性交\footnote{在中國大陸軍婚除外。}。
\item 無論是你或者人妻，都有着100\%的健康權。無論是對你還是對人妻，即使是軍婚，任何暴力都是不合法的。
\item 你和人妻有着100\%的正當防衛權利。
\end{enumerate}

這些權利，都是寫在二戰後世界各國籤署的《世界人權宣言》裏面的。所以在中國大陸以外，你也是有類似的法律保護。所以請謹記，你和人妻約會，是在行使你們戀愛自由的天賦人權。你們能享受這段快樂的約會，是包括中國共產黨在內的全世界反法西斯戰爭勝利的成果。而如果苦主有任何過激行爲，都是100\%違法甚至犯罪的行爲。有了法律和理論基礎，你就不會有任何不安的感覺，在心理上佔據高點，進而避免因爲慌張做出任何不理智的行爲。

做好心理準備之後，就應該做好應對了。筆者在美國曾經接受過專業的安全防衛訓練。本質上，面對苦主闖入，和面對入室搶劫，面對未被捕的槍手（active shooter）沒有本質區別，需要按照以下三步進行：躲藏、撤退、戰鬥。

\subsection{躲藏}

躲藏是傳統上，大家最喜聞樂見的與人妻約會的結局。我記得之前在抖音上看過幾個黃毛在窗外被鄰居拍到的影片，有幾十萬人點贊。但是實際上，很多時候躲藏並不是最佳的策略。雖然成功的躲藏和逃跑是最完美的結局，但是我們也要考慮事情的進展不是我們想象的那麼盡善盡美的話，後果也是我們可以承受的。

這裏分析一下最常見的幾個躲藏點的優劣。

首先就是窗外了。這件事要分開來看，首先如果苦主家是一樓，跳窗其實是最佳的躲藏點——因爲一跳出窗外直接就走掉了。二樓的話就要看情況了。以下是從二樓跳窗的正確步驟：
\begin{enumerate}
\item 首先要看樓層高度、地面狀況，還有你本人的身手，評估安全性。如果覺得沒有100\%的把握，不如選擇其他方式。
\item 跳窗之前手抓着窗框，盡量減少腳和地面的距離。
\item 落地之後盡量增加衝擊時間，減少衝擊力。可以通過下蹲，翻滾等方式減少衝擊力，進而減少受傷的機率。
\end{enumerate}

\textbf{如果苦主家在三樓或以上，絕對不要藏到窗外！}如果不慎墜落，可能會有生命危險。約會本來應該是快樂的事情，一切以安全第一。

另外一個很有名的隱藏地點，就是牀底。很多日本的成人電影，都有人妻的男朋友藏到牀底的情節。但是其實藏到牀底也是非常不建議的！在牀底活動空間非常有限。如果不得不長期躺在牀底，可能會造成血液循環受阻等危害身體的狀況。而且在牀底如果被發現，基本上你就成了困獸猶鬥的獸，關門打狗的狗，任人魚肉，無論是有序撤退還是反擊戰鬥都沒有辦法。所以躲藏在牀底也是非常不建議的。

比較建議的躲藏地點，是可以容納你站立的衣櫃。雖然可能躲藏效果不是很好，但是首先是一個可以靠着的站立姿勢，可以長時間保持而不至於太累。而且從站立姿勢開始，無論是撤退還是戰鬥，你都是做好準備。你可以從衣櫃的門縫觀察，聽聲音了解苦主的方位。可以做到敵在明，我在暗。一旦苦主接近，你可以立刻開門進入戰鬥狀態。無論是對後續的談判、撤退、戰鬥都是非常有利的。

一個比較接近衣櫃的藏身地點是窗簾後面。但是必須是可以觀察到周圍情況的窗簾，比如在窗簾中間的縫隙處。這樣便有和衣櫃一樣的優點。但是如果是在窗簾後但是無法觀察周圍情況的話，就有可能被敵人直接攻擊。苦主打開衣櫃門的話，你是有反應時間的。而直接攻擊窗簾後面的你，你是沒有反應時間的。

還有一個和衣櫃差不多一樣好的藏身地點，便是一些沒人的房間。如果苦主家和郭菊陽波託菲諾純水岸豪宅一樣比較大，大機率會有一些閒置的房間。如果在發現苦主回家之後能夠逃到閒置房間躲藏，也是很不錯的選擇。房間的選擇，盡量選那些比較大的，這樣如果被苦主發現，也有足夠的撤退和戰鬥的空間。

\subsection{撤退}

在躲藏之後，第二步便是撤退。

撤退分爲兩種，第一種是最好的情況，也就是在苦主發現之前離開。要能做到這一點，要在苦主注意力分散的時候從門口離開。要能夠安全離開，最好是人妻能夠轉移苦主注意力，在苦主不注意的時候從大門離開。所以能夠安全地不被發現地撤退，其實也非常考驗和人妻的信賴程度。本質上還是將自身安全交給一個女人手中。所以還是非常不建議去人妻和苦主家約會。

如果不得不面對苦主，也不需要慌張。即使你在性交過程中被發現，其實你的危險也是極低的。如果你慌張，反而容易引起風險。

面對苦主第一步，就是講明情況和給苦主普法。首先最重要的一步，就是說明你是被人妻，也就是房子的女主人邀請過來的。如果在美國，這一點尤其重要，避免被苦主裝傻開槍。跟苦主說明這一點之後，便跟他說我現在要離開了。然後開始撤離。做到這一點，一般智商平均以上的苦主，都明白你的實力，不會輕舉妄動。

撤離的時候，必須鎮定。不要慌張逃跑。正常的穿好衣服，帶好隨身物品，然後鎮定離開。手上最好拿着可以當武器的東西，比如水壺，手提電腦，有尖頭的筆，等等。然後撤離的時候，必須一直保持面對苦主，眼睛要堅定地盯着苦主的眼睛，給對方一種生理的威懾。

這種應對方式，在大部分的情況下都能安全撤離。當苦主有任何試圖物理上接近你的傾向時，必須立刻大吼：「不要碰我！」必須清晰明白地向對方表明自己的底線，就是對方沒有權利對你有任何的物理接觸。然後對對方進行警告：我是受女主人邀請來到這裏，你沒有任何權利侵犯我。任何肢體接觸都是不合法的。任何越界都會遭受到無限制的正當防衛。

一般來說到這一步，大部分智商90以上的苦主都會放棄對你進行物理攻擊。如果苦主繼續無視警告對你進行攻擊，那麼就進入第三部分，戰鬥。

當然，這個警告到戰鬥的界線的具體位置，可能因人而異。比如說，如果對方抓住你不讓你離開，這時候其實已經侵犯到你人身自由。按照法律你已經有了正當防衛權。而且當你的衣服或者手臂被抓住的時候，對後續的戰鬥，已經讓你處於很不利的地位了。所以直接進入戰鬥也是完全正當的。但是當對方沒有攻擊性行爲的時候，是否直接進入攻擊狀態，也是需要根據具體情況進行判斷。比如你是否可以掙脫，比如你對戰鬥的把握，又比如對方的精神狀態，等等。你需要綜合分析現實情況，再進行具體判斷。筆者個人來說，更傾向於對任何侵犯人身自由零容忍。不過這一方面也是因爲筆者一直對戰鬥準備足夠充分，這個界線並不一定適用於所有的讀者。

\subsection{戰鬥}

很多沒有實戰經驗的人，會以爲戰鬥只是看誰塊頭大或者力量大。而實際上，除了力量之外，士氣和謀略一樣重要。

當你決定戰鬥的一開始，你就必須明白，自己已經沒有退路。你的目標就是讓敵人喪失任何反抗力量（neutralize）。從戰鬥開始的那一刻開始，你就要明白，你不再是一個文明社會裏面生存的人類，而更像是一個在叢林中生存的野獸。此時此刻，你就是一個叢林中的雄性動物，你要和另外一隻雄性動物決鬥，爭奪交配權。這個時候沒有法律上的正確和錯誤（即使你是正當防衛），只有勝利與失敗，你一定要用你的力量和智慧，把敵人打趴下。不是你死，就是我活。明確了這一點，你便能擁有比敵人更強的士氣。

一般如果苦主對你進行攻擊，大機率已經是失去理智。這時候大機率會更傾向於蠻力而不是智慧。戰鬥的時候，必須盡量退避，閃躲，避免和對方進行近身接觸。在閃躲過苦主的蠻力攻擊之後，可以觀察苦主的弱點進行反擊。

而至於力量，要清晰地認識到，和苦主的戰鬥，並不是一個有裁判有規則的拳擊比賽或者摔跤比賽。而是一個無規則的比賽，只要能達到擊倒敵人的目標，就是勝利。所以在法律允許的最大範圍內隨身攜帶武器是必須的。在美國可以攜帶手槍，在中國可以攜帶水果刀。與其練肌肉練拳腳，不如練槍法和刀法。謹記：\textbf{寧願出現在被告席上，也不要出現在醫院ICU！}

實戰打架是一門大學問，這裏篇幅有限，僅概括性地說一些大層面的。起到拋磚引玉的作用。如果有興趣的讀者可以再自己深入研究。

\section{遇到勒索的應對}

還有一種情況，就是苦主發現之後，對你敲詐勒索金錢，否則就威脅把事情鬧大。

如果你遇到這種情況，其實應該感到開心。因爲會真正跟你拼命的，是那些心態崩潰的傳統男性。也就是那些有傳統大男子主義、力工、供養者的思維的男人。但是會敲詐你的苦主，其實心態和1998年東北大下崗騎着腳踏車把老婆送去賣淫的那些人是一樣的。他們更像是把老婆當做自己的可以利用對象和財產，如果老婆跟你睡覺能給他帶來利益，他會毫不猶豫地支持。所以這個時候，你不需要太擔心自己的人身安全問題，只需要考慮如何把經濟損失降到最低。

遇到敲詐勒索的時候，首先要評估一下，如果你和人妻約會的事情曝光，對你的影響有多大。進而評估你願意付出多少贖金。和人妻約會本身，無論是在中國和美國，都是合法的。但是如果你已婚且妻子強勢，或者你的工作或者所在的社交圈有比較高的道德要求，曝光可能會對你有很大的負面影響。比如你是一個已婚然後靠省部級老丈人提拔的體制內副處級青年幹部，你願意付的贖金自然應該比較高。但是如果你和特朗普、馬斯克和我一樣，是一個美國國籍的資產數十億的單身企業家，根本不需要怕曝光等威脅。你自然不需要付很多贖金。

關於贖金的談判，也是一門大學問。這裏篇幅有限，只能夠講一些簡單的談判方法。但是筆者這裏講的談判方法並不是唯一正確的方法。

首先一般先讓苦主自行開價。這有一個好處，就是提前知道對方的「錨定金額」。有一定機率對方一開始的出價遠低於你的底價的價格，這樣子其實你就省下一大筆錢了。

如果苦主出價低於你的心理價，忍住不要表現出開心。需要保持鎮定甚至表現出爲難的表情，無論自己有沒有錢，都說這個有點爲難。然後砍價20\%左右。快速成交然後脫身。

如果苦主出價高於你的心理價，也不需要慌張，大機率苦主也只是漫天要價碰運氣。可以直接砍價到心理價位左右。還是很有可能能夠成交的。

有一點要注意，在現場和苦主對峙的時候，因爲處於一個不熟悉的不利的戰鬥環境，可以先給或者先承諾給比較多的錢。一旦離開，其實完全可以重新談判。畢竟自由戀愛和約會並不違反任何法律，而苦主如果通過威脅獲取財物，是不合法的。在中國如果金額超過2000元人民幣，屬於數額較大的敲詐勒索罪。這時候可以再重新談判要求降低金額，甚至退還苦主已收取的財物。如果對方拒絕，甚至可以以敲詐勒索提出起訴。苦主對黃毛的索取財物的行爲定性爲敲詐勒索罪在國內已經有無數的案例。所以不需要太過在乎當場的得失，事後其實都可以重新談判。

\section{如果你是苦主}

如果你是在矛盾的另外一方的苦主，這一章會告訴你遇到妻子出軌應該怎麼做。

前面的關於黃毛面對苦主的各種反應應該如何應對，其實就是苦主的負面例子的集合。苦主首先必須明白一點：爲了一個出軌的女人和另外一個男人戰鬥，然後把自己置於危險的境地，是極其愚蠢的行爲。世界上女人幾十億，而你現在面對的，是一個大齡、已婚、沒有道德的出軌女人。現在有這個機會可以合理合法地離婚換老婆，何樂而不爲？這時候任何的搶女人、挽回都是低智商、低價值的行爲。

面對妻子出軌，一定要冷靜，然後制定計劃。目標是如何在離婚中盡量減少損失，獲取最大的利益。這件事，其實和公司解僱嚴重違規的員工是一個道理。我曾經仔細研究過紐約的自營量化基金2 Sigma起訴高亢的刑事和民事案件，2 Sigma的一整套應對讓人驚嘆。具體來說，高亢在回國創立銳天投資和衍復投資之前，在2 Sigma工作的時候，最後大約一年的時間，都在有意識的開始複製公司策略和程式碼，而且研究他不應該閱讀的公司內部商業機密文件。在公司發現高亢的這種盜竊商業機密的犯罪行爲的時候，並沒有立刻解僱他，而是放長線釣大魚，在最後幾個月，不驚動他。然後讓公司法務和IT部門合作，仔細搜集了高亢的所有違法證據。在幾個月後證據搜集得差不多，而高亢準備攜帶商業機密潛逃中國大陸的時候，和紐約法院和警方合作，迅速解僱、切斷權限、逮捕、起訴，一條龍操作，最後將他送進監獄。

2 Sigma這種專業的手法，值得每一個苦主學習。對待具體的人妻，具體的操作分爲以下四步。

第一步是情感抽離，就像一個專業、成熟的老闆，在遇到員工背叛的時候不應該對員工有個人感情一樣，一個專業的苦主，對待妻子出軌也不應該夾雜太多個人情緒。

第二步，就是在不驚動妻子的情況下，盡量多收集證據。這一步涉及很多法律專業知識，應該先諮詢律師，在律師的指導下，可以僱傭合法的第三方調查取證機構盡量收集最多的證據，以期在後續的法律鬥爭中立於不敗之地。

第三步，便是在不驚動妻子的情況下，盡量轉移所有的資產，比如可以來美國找我，把所有流動資金都跟我買成比特幣，然後放在一個除了你之外沒人知道的冷錢包裏面。然後對於其他財產，也做好所有的防禦措施防止妻子轉移資產。

第四步，便是攤牌。在攤牌的同時，應該確保所有能控制的財產控制住，然後將她趕出家門。確保她無法從物理上對你進行傷害和對家庭造成破壞。

完成了這四步之後，就可以正式進入法律程序了。

另外，如果有孩子的話，建議先進行親子鑑定。如果不幸真的有共同的子女的話，無論多生氣，還是建議保護好子女。即使大部分女人有母性，但是也看看現在北方農村那些生了小孩之後就改嫁的婦女、還有中國比生育數量還多的墮胎數量，就知道女性對子女的愛也不是必然。而且不管男人女人，出軌後不管不顧甚至仇視傷害原有子女的情況太多了。當時郭菊陽跟我說過，要是沒有小孩就好了。嚇得我以爲她要學重慶姐弟墜亡案，趕緊安撫她，叫她趕緊別亂說話也別想。所以在妻子出軌的情況下，保護好子女也是非常重要的。
